\documentclass[10pt, aspectratio=169]{beamer}
\usetheme{Madrid}
\usecolortheme{dolphin}
\usepackage{ctex}
\usepackage{xeCJK}
\usepackage{amsmath,amssymb}
\usepackage{listings}
\usepackage{tikz}
\usepackage{xcolor}
\title{408考研零基础 --- 数据结构}
\subtitle{1-15 图的基本概念与存储}
\author{}
\date{}
\setbeamertemplate{page number in head/foot}{}
\begin{document}
\frame{\titlepage}
\begin{frame}{本节目标}
\begin{enumerate}
\item 掌握图的定义和术语
\item 理解邻接矩阵和邻接表的存储结构
\item 能够根据实际需求选择合适的存储结构
\end{enumerate}
\end{frame}
\begin{frame}{图的定义与术语}
\textbf{图} $G=(V,E)$
\begin{itemize}
\item $V$ --- 顶点的非空有限集合
\item $E$ --- 边的有限集合
\end{itemize}
\vspace{0.3em}
\begin{tabular}{|c|c|}
\hline 术语 & 定义 \\ \hline
有向图 & 边有方向（弧） \\ \hline
无向图 & 边无方向 \\ \hline
度 & 与顶点相关联的边数 \\ \hline
入度/出度 & 有向图中指向/指出的边数 \\ \hline
连通分量 & 无向图中的极大连通子图 \\ \hline
强连通分量 & 有向图中的极大强连通子图 \\ \hline
\end{tabular}
\end{frame}
\begin{frame}{完全图}
\textbf{无向完全图} $|E|=n(n-1)/2$
\vspace{0.3em}
\textbf{有向完全图} $|E|=n(n-1)$
\vspace{0.5em}
\textbf{稀疏图}：边数很少\quad \textbf{稠密图}：边数接近完全图
\end{frame}
\begin{frame}{邻接矩阵}
$n$ 个顶点的 $n\times n$ 方阵，$A[i][j]=1$ 表示存在边
\vspace{0.3em}
\textbf{特点}
\begin{itemize}
\item 无向图对称，有向图不一定对称
\item 判边 $O(1)$，空间 $O(n^2)$
\item 适合稠密图
\end{itemize}
顶点 $i$ 的度：无向图 $\to$ 第 $i$ 行之和
\end{frame}
\begin{frame}{邻接表}
顶点表（数组）+ 边表（链表）
\vspace{0.3em}
\textbf{特点}
\begin{itemize}
\item 空间 $O(n+e)$，适合稀疏图
\item 判边 $O(\text{degree}(v))$
\item 遍历邻接顶点更高效
\end{itemize}
\vspace{0.3em}
\begin{tabular}{|c|c|c|}
\hline 对比 & 邻接矩阵 & 邻接表 \\ \hline
空间 & $O(n^2)$ & $O(n+e)$ \\ \hline
判边 & $O(1)$ & $O(\text{degree})$ \\ \hline
\end{tabular}
\end{frame}
\begin{frame}{本节总结}
\begin{enumerate}
\item 图分为有向图和无向图
\item 邻接矩阵空间 $O(n^2)$，适合稠密图，判边 $O(1)$
\item 邻接表空间 $O(n+e)$，适合稀疏图
\end{enumerate}
\vspace{1em}
\begin{center}
\textcolor{blue}{\textbf{下节预告：图的遍历}}
\end{center}
\end{frame}
\end{document}
